\section{Behavioural finance: biases and market behaviour}

These entries classify the systematic departures from rational decision making that matter to investors and explain how, with limits to arbitrage, they can show up in prices and trading. The investor case they refer to is \hyperref[ex:investment-planning-case]{the planning case}; manager and team behaviour is in \hyperref[meth:manager-philosophy-process-behaviour]{manager process and behaviour}.

\begin{qfdefinition}[Behavioural biases: cognitive errors compared with emotional biases; overconfidence, representativeness, conservatism, confirmation, anchoring, availability, hindsight, framing and mental accounting; loss aversion, regret and pride, self-control and ambiguity aversion; effects on decisions and remedies]\label{def:behavioural-biases}
\qffield{Statement} Rational investors in traditional finance are assumed to be risk averse, to hold rational (unbiased) expectations and to practise asset integration, judging each investment by its effect on the whole portfolio. Behavioural finance recognises investors who ``exhibit loss aversion'', ``hold biased expectations'' and ``practice asset segregation'' \citep[chap.~2, pp.~28--29]{managing-investment-portfolios-2007}. Barberis and Thaler divide the psychology into \emph{beliefs} and \emph{preferences} \citep[NBER working paper 9222 version, secs.~3.1--3.2, pp.~12--21]{wide-barberis-thaler-survey}:
\begin{itemize}
\item \emph{Overconfidence}: 98 percent confidence intervals contain the truth ``only about 60 percent of the time'', partly through self-attribution and hindsight bias (p.~12).
\item \emph{Representativeness}: probability judged by resemblance, producing base-rate neglect and sample-size neglect (the ``law of small numbers'', so four good picks make an analyst look talented) (pp.~13--14).
\item \emph{Conservatism}: too little weight on new evidence. In Edwards's urn experiment the correct posterior is 0.97 but most people answer about 0.7 (p.~14).
\item \emph{Belief perseverance and confirmation}: reluctance to seek, or to credit, contrary evidence (p.~14).
\item \emph{Anchoring}: estimates adjusted too little from an arbitrary starting value; \emph{availability}: recent and salient events weigh too heavily (pp.~14--15).
\item \emph{Prospect theory}: value defined over gains and losses, concave for gains and convex for losses, with a kink at zero (\emph{loss aversion}) and overweighting of small probabilities (pp.~16--18).
\item \emph{Framing and mental accounting}: the way problems are posed changes choices, and \emph{narrow framing} evaluates each gamble alone rather than with existing risks (pp.~19--20).
\item \emph{Ambiguity aversion}: dislike of unknown probabilities, with a ``preference for the familiar'' where people feel competent (pp.~20--21).
\end{itemize}

\emph{Cognitive against emotional.} Statman separates \emph{cognitive} shortcuts and errors (framing, hindsight, confirmation, anchoring and adjustment, representativeness, availability, confidence) from \emph{emotional} shortcuts and errors (hope and fear, happiness and sadness, anger, regret and pride, self-control, trust, mood and affect) \citep[chaps.~3--4, pp.~25--65]{behavioral-finance-second-generation-2019}. Both are shortcuts of the fast, intuitive System~1 that usually lead to good-enough choices and become errors ``when they take normal people far from their best choices''. The slow, effortful System~2 corrects them, aided by knowledge, instruments, rules and effective committees (pp.~3--6). Regret and pride are ``cognitive emotions'', the outcome of cognition and emotion together, and self-control likewise ``centers on the interaction between cognition and emotion'' (pp.~55 and~58).

\qffield{Practice} \emph{Compare and contrast.} Interpretation, drawing on the sources above. Cognitive errors arise from faulty information processing or reasoning. They are corrected by information, calculation and procedure: base rates, Bayes' rule, a checklist, an independent second opinion. Emotional errors arise from feelings about outcomes. Facts alone change them little, so advisers more often work around them with rules, defaults, staged decisions and commitment devices. Statman's list shows the border is not sharp: overconfidence has cognitive and emotional parts, and regret aversion needs foresight to operate.

\emph{Effects on decisions.} Overconfidence leads to under-diversification and excessive trading. Representativeness leads to chasing recent winners and managers with short records. Conservatism slows the response to news. Confirmation keeps theses alive after evidence turns. Anchoring ties sale decisions to purchase prices. Availability overweights vivid risks. Loss aversion and narrow framing produce reluctance to hold equities. Regret aversion and pride produce the disposition to sell winners and hold losers. Weak self-control produces overspending; excessive self-control, underspending. Familiarity produces home and employer-stock bias.

\emph{In the planning case.} Dr~A's stated tolerance for 30 percent losses and her active trading suggest overconfidence. Her attachment to the inherited stock (``my father chose it'') is emotional. The IPS answers with a quantitative risk objective based on ability, a written staged-sale rule (a self-control device) and annual review (\hyperref[ex:investment-planning-case]{planning case}). Kahneman recommends collecting each committee member's written view before discussion to ``decorrelate error'' \citep[chap.~7, ``Exaggerated Emotional Coherence (Halo Effect)'']{thinking-fast-slow-2011}, and the premortem, which ``overcomes the groupthink that affects many teams once a decision appears to have been made'' (chap.~24, ``The Premortem'').

\qffield{Conditions} The experimental findings come from laboratory and survey settings; Barberis and Thaler note that learning, expertise and incentives attenuate biases but rarely remove them (p.~15). The classifications are organising devices; neither source claims a uniform list (Statman, p.~25). Citations of Barberis and Thaler give the pages of the September 2002 NBER working paper, not of the 2003 \emph{Handbook} chapter.

\qffield{Pitfalls} Labelling every poor outcome a bias (hindsight). Treating a client's preference, such as wanting a dividend for spending discipline, as an error when it serves a real want. Assuming professionals are immune: experts can be more overconfident when feedback is limited (Barberis and Thaler, p.~15).

\qffield{See also} \hyperref[mech:biases-market-anomalies]{biases and market behaviour}; \hyperref[meth:investment-policy-statement]{IPS}; \hyperref[def:risk-aversion-utility]{risk aversion and utility}.

\end{qfdefinition}

\begin{qfmechanism}[From investor biases to market behaviour: limits to arbitrage, momentum, long-term reversal, value and post-earnings drift, closed-end fund discounts, under-diversification, excessive trading and the disposition effect]\label{mech:biases-market-anomalies}
\qffield{Statement} Behavioural finance has ``two building blocks'': limits to arbitrage and psychology. If arbitrage were unlimited, rational traders would undo mispricing caused by less rational ones; in fact arbitrage is risky and costly \citep[NBER working paper 9222 version, abstract and secs.~2.1--2.2, pp.~3--6]{wide-barberis-thaler-survey}:
\begin{itemize}
\item \emph{fundamental risk}: substitutes are imperfect, so a hedged arbitrage still loses on asset-specific news;
\item \emph{noise trader risk}: the mispricing can widen before it closes, and because arbitrageurs manage other people's money (``a separation of brains and capital''), losses can force early liquidation;
\item \emph{implementation costs}: commissions, spreads, price impact, borrowing fees and outright bans on short selling.
\end{itemize}
With arbitrage limited, the psychology of \hyperref[def:behavioural-biases]{beliefs and preferences} can affect prices. Barberis and Thaler review patterns that the CAPM does not explain (pp.~33--34): a small-firm premium; long-term reversals of three-year winners and losers (about 8 percent a year in De~Bondt and Thaler); higher returns on high book-to-market and earnings-to-price stocks; and momentum (prior six-month winner deciles beating losers by about 10 percent a year in Jegadeesh and Titman).

\emph{A belief-based account.} In the model of Barberis, Shleifer and Vishny, conservatism makes investors underreact to an earnings surprise, producing post-earnings-announcement drift and momentum. After a run of good news, representativeness makes them extrapolate and overreact, producing long-term reversals and the value effect (pp.~38--39).

\emph{Closed-end funds.} Funds typically trade at about a 10 percent discount to NAV that varies over time, start at a premium and converge to NAV at termination. Rational explanations account for parts of this but not all of it, and fluctuating investor sentiment with limited arbitrage is offered as an explanation (sec.~6.1, p.~44).

\emph{Investor behaviour.} Home bias (French and Poterba: 94, 98 and 82 percent of US, Japanese and UK equity held domestically), employer stock and naive $1/n$ diversification fit ambiguity aversion and familiarity. Excessive trading lowers net returns, most for those who trade most, consistent with overconfidence. Investors sell winners rather than losers, the disposition effect, though tax and subsequent returns favour the opposite (pp.~47--50).

\qffield{Conditions} The anomaly figures come from the samples Barberis and Thaler cite, mostly US data before 2000, and the paper notes that the size premium fell in the following decade (n.~23, p.~33). Rational risk-based explanations compete with behavioural ones for several patterns, and a pattern's existence does not identify its cause. That anomalies persist only where arbitrage is limited is the theory's prediction, not a proof that every pattern is mispricing.

\qffield{Practice} Interpretation: limits to arbitrage tell an investor when a strategy that exploits a bias is likely to hurt before it pays (noise trader risk), and why capacity and shorting costs matter. For an adviser, the investor-behaviour findings explain why clients' realised returns trail fund returns and why the IPS needs rules.

\qffield{Pitfalls} Reading a published backtest of an anomaly as evidence of exploitable mispricing after costs and after its publication. Assuming that because individual biases exist, prices must be wrong; aggregation and arbitrage can offset them. Ignoring that the same arbitrageurs can be forced sellers in a crisis.

\qffield{See also} \hyperref[def:behavioural-biases]{behavioural biases}; \hyperref[def:alpha-model-taxonomy]{alpha model taxonomy}; \hyperref[meth:skill-luck-significance-tests]{skill against luck}.

\end{qfmechanism}
